\section{Discussion in full}
\label{sec:discussion-full}
\subsection{What a union-recall number means}

A union-recall figure depends on the false-positive rate it cost, so recall quoted without that
rate cannot be carried to another setting. We make this as a methodological point. The measured
support we tried to give it came from a second substrate that was run three times and closed
without a quotable number (\S\ref{sec:inreview}). The point needs no second substrate:
recall and false positives move together with any severity threshold, and a critic can be tuned
to any recall by blocking more. Our records support a narrower version on the substrate we
report. The same auditor at two severity rules sits at two different operating points
(\S\ref{sec:auditor}), and the shipped auditor's 30.0\% [20.0, 40.7] recall is quoted only beside
the 16.0\% [10.1, 22.3] it cost.

The cleaner experiment for the general claim is a severity sweep within one substrate. Our sweep
re-grades archived findings at four decision rules and shows each family's recall and
false-positive rate moving with the rule (Figure~\ref{fig:sweep}). Its rules are deterministic
and do not place the families at a common operating point. A randomised mixture of two archived
rules could reach one in expectation without new model calls, but it would be a new post-hoc
policy needing its own uncertainty analysis, and we have not built it. Two substrates differ in
specification length, solution length, problem population, defect distribution and suite
construction all at once, so comparing them could never have isolated the operating point.

\subsection{Repeated reading buys nothing where the verdict does not vary}

Union-of-$K$ is the instrument this paper uses to look for a ceiling, and part of what it cannot
measure follows from arithmetic alone. A union over readings that do not differ cannot grow, so
the curve is flat for any route that returns the same verdict however often it is asked.

Our records establish the design fact that makes this matter here. The same-vendor routes are sent
temperature $0$, while the cross-vendor route is sent no sampling parameter
(\S\ref{sec:auditor}). From $K = 1$ to $8$ the first family's union gains 3.0 points and the
second's gains 19.3. No temperature-matched arm was run, so we cannot say how much of that gap is
sampling and how much is the operating point. We treat it as a confound on our own contrasts.

For best-of-K, union-of-K or panel aggregation over critics, what follows is the arithmetic: a
route whose verdict does not vary between readings gains nothing by being read again. We measured
per-instance flag counts, not any distance to a decision boundary, and ran no panel aggregation,
so we offer no law for how much a varying route gains. The diminishing shape of the curves is not
itself evidence of a ceiling either. Exact subset averaging makes each reading's marginal gain no
larger than the last by construction, and the curves measure the size of those gains.

\subsection{What the ceiling is made of}

We do not rank these contributions, because they were measured on different populations, depths
and rules. Two of them share a scale. Eight readings of one auditor buy 19.3 points of union
recall. Twenty readings across three families buy a pooled 48.2\% [36.7, 60.0] at 36.0\% false
positives on correct code, and leave the complement, 51.8\% [40.0, 63.3] of the defect
population, untouched. Changing which model reads moves the same quantity (union recall at $K=8$
over the same 110 instances) by 26.4 points, more than repetition's 19.3. A same-size same-vendor
substitution moved it by only 12.7, and \S\ref{sec:auditor} gives three reasons not to read either
number as auditing ability. Changing one sentence of the rulebook moves flags by 26.8 points
[13.6, 40.4], an exploratory contrast on 56 instances at one reading, at $+12.5$ points on correct
code ([$-3.6$, $+26.8$]) and with no established effect on hidden-suite outcomes; it is not on the
scale of the other two.

A fourth contribution is not the auditor's at all, and we can describe it only post hoc. Under a
rule that asks whether the specification's prose determines the hidden suite's expected value, 44
of the 57 never-flagged defects fail. We tried to build a population where the prose does
determine the answer. Review established that no filter in it enforces that property, so it
licenses description only (\S\ref{sec:inreview}).

The association between never being flagged and never being determined is empirical. Selecting
the missed instances does not fix their labels, and against caught instances the third rater's
post-hoc contrast is $+23.3$ points $[+2.0, +43.6]$. That rater agreed with itself on only 7 of
11 repeated items, and the interval is conditional on its recorded labels. Selection and
confounding rule out a causal reading, so the contrast locates the misses and does not estimate
how much the benchmark contributes.

One intervention changes the specification instead of conditioning on it
(\S\ref{sec:inreview}). Adding a behavioural rule to the prose moved diagnosis on 9 of 32 such
instances, all of them where the added rule agreed with the hidden suite and none of the 23 where
it contradicted it. This is the only evidence here that manipulates the specification. Its
population was selected for having been missed and called undetermined, and it cannot separate a
specification that settles the question from one that settles it correctly.

The constructed population has two further limitations, one measured and one found in review. A
frontier model separates its edits from natural ones at 96.7\% on the items it answered, and
between 79.9\% and 97.3\% once the items it would not classify are counted at either extreme. It
does so \emph{by a feature it does not identify}: the two arms share no problem, and the model is
one of the two that gated the population. And because the construction does not enforce
specification-determinedness, it licenses \emph{no} bound on the measured ceiling.

\subsection{Flagging is not naming}

An audit that blocks an increment for the wrong reason has still blocked it, and a recall number
counts that as a catch. For one reading of the shipped auditor the finding texts were kept and
adjudicated blind by two raters. The auditor returned a finding on 9 of 110 defective instances
and asserted the actual failing behaviour on 5 to 7 of them, depending on how the two raters are
combined. The raters agreed at $\kappa$ = 0.391. We report that beside the rate because a moderate
agreement is part of what the rate is worth. That study's preregistered kill condition fired on
this result, and \S\ref{sec:inreview} states its headline in the form the kill requires. The
figure describes one reading. It is not a naming rate for the eight-reading union, which would
need the texts of all eight readings, and those were not archived.

\subsection{What the science studies add}

They give the science profile's design rule (the model names
evidence, code verifies it) an empirical reading and a boundary. On results, number$\to$source
flagged every inconsistency between a report, its results file and its log, and could not see a
result that was wrong consistently; the model flagged 23 of 25 such fabrications, 21 with a sound
finding about the fault (post hoc). On data, a validator written from the card flagged what a rule can
state and missed what needs reading across columns and units; the model flagged 18 of the 28 such
items the validator missed with a finding about the fault (post hoc).
Where both tiers ran, the deterministic one flagged no clean item and the model carried the errors;
on data the union was better than either part, and on results the model's flags already contained
the profile's, whose findings its prompt carried. Neither study shows that the model tier scales: the
artefacts were small, the model's fabrication findings are hand calculations (several wrong), and it
missed every fault that needs comparison across many rows. A model reading is a useful second tier
where the artefact fits in the reader's working memory, not a substitute for executing the code or
computing over the data.

\textbf{Clean items must be clean by the auditor's standard.} In the results and data studies the
auditor flagged defects in our own construction: a card and file that named a column differently, a
file outside the audited scope, a report claiming a run no file performed, a provenance sentence false
for inputs passed as literals, and a unit label contradicting the code's documentation. Each was a
correct flag on an item we had built as clean. We caught two in pilots, one when every verdict came
back as an escalation, and the rest by reading every flag on a clean item, which the registrations
required. In the code study the task we wrote named the whole problem when one step was delivered,
and a sample of the auditor's flags on correct code objected to exactly that; there we read a sample,
not every flag. Without such reading, a false-positive rate can measure the evaluator.

\subsection{Limits}
\label{sec:limits}

\begin{enumerate}\itemsep1pt
\item \textbf{Sampling differs between families.} One route is sent temperature $0$ and the other
  no sampling parameter, which bounds how far one union can grow before any difference in ability
  is measured. Every cross-family union-of-K contrast on general-purpose code inherits this,
  and it is the confound we would fix first. In the science studies the same-vendor route runs at
  its default sampling through a command-line route, a different asymmetry.
\item \textbf{No matched operating point.} No auditor is compared with another at a matched
  false-positive rate. Our sweep re-grades archived findings at four deterministic decision
  rules, under which one family has no point between 3.3\% and 34.0\% false positives. We built
  no randomised policy between them, so the paper cannot say whether the families differ at the
  same operating point.
\item \textbf{Interval coverage.} The primary general-code intervals are problem-cluster
  percentile bootstraps over 56 clusters (the science studies' have their own cluster counts, and the
  data study resamples items within dataset). A simulation of an analogous single-rate interval (600 resamples
  rather than 10{,}000, tested rates only) covers near nominally at the headlines' rates and
  anticonservatively below about $0.10$ (\S\ref{sec:stats}). For the paired contrasts and
  conditional shares coverage is unmeasured, and no sample size was justified in advance by a
  detectable-effect calculation.
\item \textbf{Substrate.} The substrate is short, self-contained Python functions with executable
  hidden tests, so nothing here measures auditing a change to a large program with no oracle.
  Both benchmarks are old enough that memorisation cannot be ruled out.
\item \textbf{One generator.} The defect population is one model's failure mode.
\item \textbf{Raters.} The residual rubric's raters were two models, one of them the agent that
  wrote the rubric, which was not blind to the earlier labels. One external label set agrees with
  them across eleven overlapping tasks.
\item \textbf{The constructed population.} It did not enforce specification-determinedness, and
  its edits are separable from natural code at 96.7\% of the items the probe answered
  (79.9--97.3\% once unanswered items are counted at either extreme) by something the probe does
  not identify.
\end{enumerate}
Each limit is recorded in the study that carries it, together with the number it would change.

\subsection{What we would do next}

The sweep re-grades archived findings under four rules and shows that the auditor contrast
depends on the rule (\S\ref{sec:auditor}). Re-grading only moves a cut through findings already
produced, so it cannot bring every family to a matched operating point: one family has no point
under the four deterministic rules between 3.3\% and 34.0\% false positives. The experiment this
paper most needs is to re-ask, instructing each family at graded severity and reading afresh,
which can place the families at a common false-positive rate and decide whether the auditor
contrast survives there. Matching the sampling parameters across families would remove the
confound above, and reading the cross-vendor family past eight times would turn a statement about
our budget into one about a curve.

We ran a first version of the experiment we most wanted, holding the code and the visible tests
fixed and varying only the specification (\S\ref{sec:inreview}). Two gaps remain, and they decide
what it means. The clarifier inferred the missing rule from failing inputs and was wrong on most
instances, so the rule's correctness was measured, not manipulated. Writing the rule from the
oracle for some instances, and a deliberately wrong rule of matched length for others, would
separate a specification that settles the question from one that settles it correctly. Its
manipulation check also did not show the placebo inert, and a placebo rated no more determining
than the original would. Both are cheap, because the code, the tests and the adjudication rubric
already exist.

Beyond those: adjudicating more than one reading of the shipped auditor would let a naming rate be
compared with a union flag rate at the same K. A second generator would show whether the
residual's character belongs to the benchmark or to one model's failures. A population of natural
defects independently known to be specification-determined would separate the injected
population's artificiality from the property it was built to isolate. And a substrate whose ground
truth is not a hidden test suite would test the assumption every ceiling here rests on, that a
test suite asked the right question.
